\section{Introduction}\label{sec:introduction}

The Boltzmann equation describes the transport and binary collisions of a
dilute gas. Its collision law conserves mass, momentum, and kinetic energy
at each spatial point, while its entropy decreases. A weak solution theory
must retain enough of these properties to distinguish the equation from
a balance law with an unidentified collision defect. This paper constructs
two global solutions with the same initial datum for which all five local
collision-invariant balances hold exactly.

The connection between binary collisions and entropy dissipation goes
back to Boltzmann's kinetic theory~\cite{Boltzmann1872}. For the
spatially inhomogeneous Cauchy problem, DiPerna and Lions established
the large-data theory of global existence and weak stability for
renormalized solutions~\cite{DiPernaLions1989}. On periodic space the
initial assumptions are finite mass, energy, and entropy, and the
solutions obey the entropy-production inequality
\cite[Theorem~1 and Remark~2]{DiPernaLions1991}.
Renormalization addresses a basic difficulty: these bounds alone need
not make the quadratic collision term locally integrable.

Exact conservation is an additional issue at this level of regularity.
The periodic formulation in
\cite[Theorem~4.1]{LevermoreMasmoudi2010} gives local mass conservation,
but allows a nonnegative matrix-valued defect in the momentum balance
and a corresponding defect in total energy. It does not automatically
give exact local momentum and energy conservation. The present result
concerns solutions that obey both of these local identities, with the
energy flux integrable and with the prescribed initial traces.
Their separate collision gains and losses are integrable with every
polynomial velocity weight on each finite time interval.

Uniqueness is available under stronger control. Kaniel and Shinbrot
proved local existence and uniqueness in classes bounded by a multiple
of a Maxwellian~\cite{KanielShinbrot1978}. Ukai obtained global
existence and uniqueness for small periodic perturbations of a Maxwellian
in weighted spatial Sobolev spaces~\cite[Theorem~1.2]{Ukai1974}.
Duan, Huang, Wang, and Yang later allowed large but bounded amplitudes
on the torus or in the whole space, assuming small relative entropy and
a small spatially integrated velocity supremum of the
perturbation~\cite[Theorem~1.1]{DuanHuangWangYang2017}.
Our datum instead has unbounded amplitudes at shrinking spatial scales.
The small relative entropy used in its construction therefore does not
place it in these uniqueness classes.

For the spatially homogeneous hard-sphere equation, finite mass and
energy suffice for uniqueness within the conservative class
\cite[Theorem~1.1]{MischlerWennberg1999}. Wennberg's homogeneous
nonuniqueness examples for hard potentials allow increasing
energy~\cite{Wennberg1999}. The mechanism studied here uses spatial
concentration and preserves energy locally. Thus the distinction between
homogeneous and inhomogeneous evolution, as well as the precise
conservation requirements, matters for comparison.

Recent instability results address a different aspect of the Cauchy
problem. Chen and Holmer prove failure of uniform continuity in weighted
Sobolev spaces for the constant collision kernel on
\(\R^3_x\times\R^3_v\)~\cite[Theorem~1.3]{ChenHolmer2024}.
Chen, Shen, and Zhang extended this analysis to a range of cutoff soft
potentials~\cite{ChenShenZhang2024}. For hard spheres on the whole space,
Chen, Guo, Shen, and Zhang prove failure of uniform continuity from
Gaussian-weighted initial data to the unweighted \(L^2_vH^s_x\) norm
for \(0\le s<1\)
\cite[Theorem~1.1]{ChenGuoShenZhang2026}. These results compare nearby
initial data, whereas the two solutions constructed here have exactly
the same datum. For the Krieger--Strain (isotropic Landau) model,
Gismondi, Golding, and Novack construct rough stationary solutions whose
mollified collision operators tend to zero, and conjecture nonuniqueness
of renormalized solutions to the inhomogeneous Landau and Boltzmann
Cauchy problems
\cite[Section~1.2]{GismondiGoldingNovack2026}.

\subsection{The equation and the solution class}

Let \(\T^3=\R^3/\mathbb Z^3\) have volume one and put
\(\Omega=\T^3_x\times\R^3_v\). Write
\(\D=\partial_t+v\cdot\nabla_x\). For \(v,v_*\in\R^3\) and
\(\omega\in S^2\), define the elastic collision by
\[
 a=(v-v_*)\cdot\omega,\qquad
 v'=v-a\omega,\qquad v_*'=v_*+a\omega.
\]
The sphere carries ordinary surface measure. The hard-sphere equation is
\begin{equation}\label{eq:collision}
 \D F=Q(F,F)=Q^+(F,F)-Q^-(F,F),\qquad F\ge0,
\end{equation}
where
\[
\begin{split}
 Q^+(F,F)(v)&=\int_{\R^3\times S^2}|a|F(v')F(v_*')\dd v_*\dd\omega,\\
 Q^-(F,F)(v)&=\int_{\R^3\times S^2}|a|F(v)F(v_*)\dd v_*\dd\omega.
\end{split}
\]
Every factor has the same time and position. We use primes and stars with
the same convention in the entropy dissipation below.

Set \(M(v)=(2\pi)^{-3/2}e^{-|v|^2/2}\) and \(G=F/M\). Define
\begin{equation}\label{eq:entropy}
\begin{split}
 \mathcal H_M(F)&=\int_\Omega
       \bigl(F\log(F/M)-F+M\bigr)\dd x\dd v,\\
 \mathcal D(F)&=\frac14\int_{\T^3\times\R^3\times\R^3\times S^2}
 |a|(F'F_*'-FF_*)\log\frac{F'F_*'}{FF_*}
       \dd x\dd v\dd v_*\dd\omega.
\end{split}
\end{equation}
The dissipation integrand takes its nonnegative lower-semicontinuous
extended value at zero. We set \(0|\log0|=0\).

\begin{definition}[Entropy solution with local conservation]
\label{def:admissible}
Let \(F_0\ge0\) have finite mass, energy, and absolute entropy.
A global entropy solution with local conservation is a nonnegative
function \(F\) satisfying the following conditions.
\begin{enumerate}[label=\textup{(\roman*)}]
\item \(F\in C_w([0,\infty);L^1(\Omega))\), \(F(0)=F_0\), with
continuity against every bounded measurable test. For every finite \(T\),
\[
 \sup_{0\le t\le T}\int_\Omega
 F(t)(1+|v|^2+|\log F(t)|)\dd x\dd v<\infty.
\]
\item The two functions \(Q^\pm(F,F)/(1+G)\) belong to
\(L^1([0,T]\times\Omega)\) for every finite \(T\).
For every \(\beta\in C^1([0,\infty))\) with \(\beta(0)=0\) and
\(|\beta'(z)|\le C_\beta/(1+z)\), and every smooth \(\varphi\) periodic
in \(x\) and compactly supported in \([0,\infty)\times\R^3_v\),
\begin{equation}\label{eq:renormalized}
\begin{split}
 &\int_0^\infty\!\!\int_\Omega
 \left[M\beta(G)\D\varphi+\beta'(G)Q(F,F)\varphi\right]
                  \dd x\dd v\dd t\\
 &\hspace{15mm}
 +\int_\Omega M\beta(F_0/M)\varphi(0,x,v)\dd x\dd v=0.
\end{split}
\end{equation}
\item Total momentum is conserved. At every \(t\ge0\),
\[
 \int_\Omega |v|^2F(t)\dd x\dd v\le\int_\Omega |v|^2F_0\dd x\dd v,
 \qquad
 \mathcal H_M(F(t))+\int_0^t\mathcal D(F(s))\dd s
      \le\mathcal H_M(F_0).
\]
\item For every finite \(T\),
\[
 \int_0^T\!\!\int_\Omega (1+|v|^3)F\dd x\dd v\dd t<\infty.
\]
For every \(\psi\in\{1,v_1,v_2,v_3,|v|^2/2\}\) and
\(\phi\in C_c^\infty([0,\infty)\times\T^3)\),
\begin{equation}\label{eq:local-balances}
 \int_0^\infty\!\!\int_\Omega F\psi\D\phi\dd x\dd v\dd t
   +\int_\Omega F_0\psi\phi(0,x)\dd x\dd v=0.
\end{equation}
\end{enumerate}
\end{definition}

The last condition is local conservation of mass, momentum, and kinetic
energy, including the initial traces. In particular it is stronger than
requiring only the corresponding global inequalities.

\begin{theorem}\label{thm:main}
There exist \(R<\infty\) and a nonnegative datum \(F_0\) with
\[
 \supp_vF_0\subset B_R(0),\qquad
 \int_\Omega F_0(1+|v|^2+|\log F_0|)\dd x\dd v<\infty,
\]
for which \eqref{eq:collision} has two global entropy solutions with local
conservation in the sense of Definition~\ref{def:admissible}, differing on
a set of positive \((t,x,v)\)-measure. Both solutions belong to
\(C([0,\infty);L^1(\Omega))\). Moreover, for every finite \(T\) and
every \(k\ge0\), they satisfy
\begin{equation}\label{eq:strong-collision-integrability}
 \int_0^T\!\!\int_\Omega
 \langle v\rangle^k\bigl[Q^+(F,F)+Q^-(F,F)\bigr]
                       \dd x\dd v\dd t<\infty,
 \qquad \langle v\rangle=(1+|v|^2)^{1/2}.
\end{equation}
\end{theorem}

\subsection{The construction and its new global estimate}

The starting point is the concentrating-jet mechanism of
\cite{OpenAIPeriodic}. That construction produces nonuniqueness in an
entropy class. Its continuation theorem does not supply the exact local
momentum and energy identities required here. We retain the jet mechanism
and construct the global continuation directly. The kernel estimates,
the detection of a nonzero rescaled solution, and its branching growth
are reproduced below with attribution; the new estimates concern the
interaction of those jets with a surrounding gas over arbitrarily long
times.

A cold jet at scale \(s\) occupies a thin transverse annulus and has a
small velocity width. Its velocity-integrated density can be of order \(1/s\), but its total
mass is of order \(s^2\). Geometrically decreasing scales therefore allow
singular collision frequencies near time zero while retaining finite
mass, energy, and entropy. The opposite axial streams are separated into
fine spatial stripes. Under dilation those stripes average to a
two-stream medium; a small additional population then follows a linear
equation in that medium.

There are two obstacles to extending this mechanism while preserving all
local conservation laws. First, thin jets eventually meet their periodic
copies. Second, an entropy bound alone does not give the separate
integrability of gain and loss needed to pass every local balance to a
limit. We address both by giving each scale a finite lifetime as a
separate component. The jet axis is chosen in a quadratic-irrational
direction, which prevents intersections before that lifetime ends. A
nearly Maxwellian bath supplies a uniform positive collision frequency
after a fixed short time. The jets then decay exponentially, so that each
can be added to a remainder with a uniform weighted bound at the end of
its lifetime.

The central continuation estimate combines three facts. On a fixed
weighted order interval, small velocity mass makes the normalized gain
small. Small relative entropy controls the spatial integral of the
excess over the Maxwellian. Finally, integration over velocities along
free flights turns this spatial smallness into a pointwise mass bound.
These facts close a global weighted supremum estimate even though the
initial cold components have no common finite supremum. The same
construction makes the contributions of the small scales to the
weighted collision integrals summable. This yields
\eqref{eq:strong-collision-integrability}, which is the decisive input
for exact local conservation and strong time traces.

For nonuniqueness, the approximation with no additional initial
population will be called dormant. In a second family, a vanishing seed
is varied until its evolution attains a prescribed time-dependent bound.
Dilation at an attaining time produces a nonzero solution of the linear
equation in the two-stream medium; nonconcentration prevents the attained
size from disappearing. Branching growth forces the rescaled limit to
have a finite time horizon. The dormant solution, meanwhile, tends to
zero away from the velocity axis under the normalized dilation by the
original horizon. Finiteness makes the two spatial and temporal dilation
scales comparable. If the two global limits agreed at arbitrarily small
horizons, comparing their dilations would force the nonzero limit to
vanish.

\paragraph{Organization.}
Section~\ref{sec:framework} fixes the gain normalization and the velocity
weight. Section~\ref{sec:kernel} proves the kernel estimates.
Sections~\ref{sec:construction}--\ref{sec:early} construct the data,
control the cold jets, and select the short-time approximations.
Section~\ref{sec:global} proves the global continuation estimate, and
Section~\ref{sec:compactness} proves admissibility and local conservation
of the limits. Sections~\ref{sec:dilation}--\ref{sec:branching} establish
the nonzero rescaled limit and its finite horizon.
Section~\ref{sec:separation} completes the proof of
Theorem~\ref{thm:main}.
